% !TeX root = ../../Book4.tex
% 纲要标题：## 附录 D　接触拓扑与填充障碍
% 纲要位置：几何学计划纲要.md，第 4664 行。
% 正文状态：未撰写；本文件目前是框架。

% 纲要内容开始
% 
% > 保留三维及高维两条支线，各自说明维数、紧性和可用不变量。
% 
% * **Legendrian 的整体不变量**
%   * 前图及尖点
%   * Thurston–Bennequin 与旋转数
%   * 稳定化和同位的区别
% 
% * **填充和配边**
%   * 强弱精确填充
%   * 典型不可填充模型及引用定理
%   * 伪全纯曲线产生的障碍
% 
% * **高维柔性与刚性**
%   * 形式解与实际解
%   * loose Legendrian 的准确对象
%   * h-原理的条件和适用界线
% 
% ---
% 
% 纲要内容结束

\chapter{接触拓扑与填充障碍}
\label{g4:app:D}
\BookChapterNumber{g4:app:D}

待撰写

\section{Legendrian 的整体不变量}
\label{g4:sec:D01}
待撰写

\subsection{前图及尖点}
\label{g4:subsec:D01:01}
待撰写

\subsection{Thurston–Bennequin 与旋转数}
\label{g4:subsec:D01:02}
待撰写

\subsection{稳定化和同位的区别}
\label{g4:subsec:D01:03}
待撰写

\section{填充和配边}
\label{g4:sec:D02}
待撰写

\subsection{强弱精确填充}
\label{g4:subsec:D02:01}
待撰写

\subsection{典型不可填充模型及引用定理}
\label{g4:subsec:D02:02}
待撰写

\subsection{伪全纯曲线产生的障碍}
\label{g4:subsec:D02:03}
待撰写

\section{高维柔性与刚性}
\label{g4:sec:D03}
待撰写

\subsection{形式解与实际解}
\label{g4:subsec:D03:01}
待撰写

\subsection{loose Legendrian 的准确对象}
\label{g4:subsec:D03:02}
待撰写

\subsection{h-原理的条件和适用界线}
\label{g4:subsec:D03:03}
待撰写
